\documentclass[11pt,a4paper]{article}

% Packages
\usepackage[utf8]{inputenc}
\usepackage[T1]{fontenc}
\usepackage[margin=0.75in]{geometry}
\usepackage{enumitem}
\usepackage{hyperref}
\usepackage{titlesec}
\usepackage{xcolor}

% Formatting
\pagestyle{empty}
\setlength{\parindent}{0pt}
\setlist[itemize]{leftmargin=*, noitemsep, topsep=3pt}

% Section formatting
\titleformat{\section}{\Large\bfseries}{}{0em}{}[\titlerule]
\titlespacing*{\section}{0pt}{12pt}{6pt}

\titleformat{\subsection}[runin]{\bfseries}{}{0em}{}
\titlespacing*{\subsection}{0pt}{8pt}{0pt}

% Colors
\definecolor{linkcolor}{RGB}{0,102,204}
\hypersetup{
    colorlinks=true,
    urlcolor=linkcolor
}

\begin{document}

% Header
\begin{center}
    {\LARGE \textbf{ERIC MAYEUX}}\\
    \vspace{4pt}
    {\large Gestionnaire de produit GIA | IAM Product Manager}\\
    \vspace{4pt}
    Montréal, QC, Canada\\
    \href{mailto:mayeux.e@gmail.com}{mayeux.e@gmail.com} | 438-884-7645
\end{center}

\vspace{8pt}

\section*{Profil Professionnel}
\noindent
\textbf{Gestionnaire de produit} avec 10+ ans d'expérience en \textbf{transformation numérique} et livraison Agile dans le secteur bancaire. Expertise en \textbf{gestion du backlog}, \textbf{cérémonies Agile} (Sprint planning, raffinement, démos, rétrospectives) et \textbf{validation des livrables}. Fort sens de la \textbf{cybersécurité} et de la conformité applicative (Zero Trust, gouvernance des accès). Capacité à servir de pont entre équipes métiers, sécurité et techniques. Certifié \textbf{PMP} (en cours), CSM, AWS Cloud Practitioner.

\section*{Compétences Techniques}

\textbf{Vision produit \& stratégie:} Feuille de route (roadmap), priorisation, initiatives de sécurité et conformité, vision produit GIA

\textbf{Gestion du backlog \& livraison Agile:} User Stories, exigences fonctionnelles, priorisation backlog, Sprint planning, raffinement, démos, rétrospectives, acceptance criteria

\textbf{Gouvernance \& Cybersécurité:} Gestion des accès, conformité, Zero Trust, RBAC, audit, réglementation, séparation des tâches (SoD)

\textbf{Collaboration \& Coordination:} Interface équipes métiers/TI/sécurité, gestion des parties prenantes, facilitation des décisions d'architecture

\textbf{Outils:} Jira, Confluence, Datadog, Splunk, tableaux de bord KPI, Azure DevOps

\textbf{Méthodologies:} Agile, Scrum, SAFe, Kanban, TDD, BDD, PMP

\textbf{Intelligence Artificielle:} Agents IA, automatisation, Playwright MCP, Crew AI, RPA

\vspace{6pt}

\section*{Réalisations}
\begin{itemize}
    \item Pilotage de la \textbf{vision produit} et de la feuille de route qualité à la BNC (contexte bancaire réglementé)
    \item \textbf{Priorisation et gestion du backlog} en alignement avec les enjeux de conformité et de sécurité
    \item Participation aux \textbf{cérémonies Agile} (Sprint planning, raffinement, démos, rétrospectives)
    \item Maîtrise des enjeux de \textbf{cybersécurité} et de gouvernance dans un contexte bancaire (BNC)
    \item Coordination entre équipes métiers, TI et sécurité pour assurer la cohérence des livrables
    \item \textbf{Amélioration continue} des processus GIA et mise en place de métriques de conformité
\end{itemize}

\section*{Expérience Professionnelle}

\subsection*{Scrum Master -- Gestion de projet | Project Manager} \hfill \textit{Novembre 2025 -- Présent}\\
\textbf{Banque Nationale du Canada (BNC)} -- Montréal, QC\\
Pilote la \textbf{feuille de route produit} dans un contexte de \textbf{transformation numérique} bancaire.
\begin{itemize}
    \item Recueille, analyse et challenge les \textbf{besoins métiers} et de sécurité ; traduit en User Stories
    \item \textbf{Priorise et gère le backlog} produit en lien avec les objectifs de conformité et de gouvernance
    \item Anime les \textbf{cérémonies Agile} : Sprint planning, raffinement, démos, rétrospectives
    \item Valide les livrables selon les \textbf{acceptance criteria} ; assure la qualité fonctionnelle
    \item Suit les \textbf{indicateurs de performance} produit (adoption, conformité, délais) via tableaux de bord
    \item Interface entre équipes TI, métiers et \textbf{cybersécurité} ; facilite les décisions d'architecture
\end{itemize}

\subsection*{Intégrateur Sénior Qualité | Senior Quality Integrator} \hfill \textit{Octobre 2023 -- Novembre 2025}\\
\textbf{Banque Nationale du Canada (BNC)} -- Montréal, QC\\
\begin{itemize}
    \item Définit et priorise les \textbf{exigences fonctionnelles} de sécurité et qualité pour le train
    \item Gestion des environnements, conformité, traçabilité des données via Datadog, Splunk, Jira
    \item \textbf{Amélioration continue} des pratiques GIA et sécurité applicative
\end{itemize}

\subsection*{Intégrateur Qualité -- SDET | Quality Integrator -- Software Development Engineer in Test} \hfill \textit{Janvier 2022 -- Octobre 2023}\\
\textbf{Banque Nationale du Canada (BNC)} via Groupe Astek -- Montréal, QC\\
\begin{itemize}
    \item Analyse statique de sécurité, tests de charge, \textbf{conformité} et gouvernance qualité
    \item Stack : Selenium, AppliTools, RestAssured, JMeter, Gatling, K6
\end{itemize}

\subsection*{QA Automaticien | QA Automation Engineer} \hfill \textit{Juin 2021 -- Décembre 2021}\\
\textbf{AODocs} -- Paris, France\\
\begin{itemize}
    \item Automatisation tests non-régression, POC Flutter ; \textbf{acceptance criteria} et rapports Allure
\end{itemize}

\subsection*{Automaticien CI/CD | DevOps Automation Engineer} \hfill \textit{Janvier 2020 -- Juin 2021}\\
\textbf{ENGIE DIGITAL} via Groupe Hénix -- Paris, France\\
\begin{itemize}
    \item Administration Jira, KPIs, gouvernance projets ; scripting Python/Shell
\end{itemize}

\subsection*{Product Owner} \hfill \textit{Juin 2019 -- Septembre 2019}\\
\textbf{Hénix} -- Montrouge, France\\
\begin{itemize}
    \item \textbf{Gestion du backlog}, priorisation, recette ; certification Jira ACP-600
\end{itemize}

\subsection*{Développeur Logiciel | Software Developer} \hfill \textit{Mai 2017 -- Mai 2019}\\
\textbf{MEDIAPOST} via Groupe Hénix -- Paris, France\\
\begin{itemize}
    \item Développement back office, webservices SOAP/REST ; \textbf{conformité} et traçabilité des données
\end{itemize}

\subsection*{Consultant QA \& Automatisation | QA Automation Consultant} \hfill \textit{Avril 2016 -- Avril 2017}\\
\textbf{ENGIE IT} via Groupe Hénix -- Saint-Ouen, France\\
\begin{itemize}
    \item Gouvernance opérationnelle, gestion des accès applicatifs, administration HP ALM
\end{itemize}

\section*{Certifications \& Formations}

\textbf{PMP Certification} (en cours) \hfill \textit{2026}

Project Management Professional -- Project Management Institute (PMI)

\textbf{AWS Cloud Practitioner} \hfill \textit{2026}

Amazon Web Services (AWS)

\textbf{Certified Scrum Master (CSM)} \hfill \textit{2024}

Scrum Alliance

\textbf{Jira ACP-600 -- Project Administration in Jira Server} \hfill \textit{2019}

Atlassian Certified Professional

\textbf{Cursus Automaticien et DevOps} \hfill \textit{2019}\\
\textit{AFCEPF -- Montrouge, France}

\textbf{Cursus Architecture logiciel \& Analyste informaticien - Certifié Master II} \hfill \textit{2017}\\
\textit{AFCEPF-HENIX -- Malakoff, France}

\section*{Formation Académique}

\textbf{Licence en Gestion (Bachelor's Degree in Management)} \hfill \textit{2011}\\
École Supérieure de Gestion (E.S.G) -- Paris, France

\section*{Compétences Linguistiques}

\textbf{Français:} Langue maternelle (Native)\\
\textbf{Anglais:} Niveau Avancé (Advanced / Professional Working Proficiency)

\end{document}
